\section{Grundlagen}
\subsection{Geschäftsprozesse}
% Plan 2.1 (Stand 28.10.2026), wird Satz für Satz abgesegnet:
%  1. Was ist ein Geschäftsprozess (Definition mit Quelle)
%  2. Prozesse werden hierarchisch zerlegt: Prozess -> Teilprozess -> Task
%  3. Formale Grundlage BPMN 2.0: Task = atomare Aktivität, Teilprozess = zusammengesetzte Aktivität
%  4. Zahl der Ebenen ist Konvention: APQC PCF nutzt fünf Ebenen
%  5. Warum die Ebenen für Compliance wichtig sind: Vorschriften gelten unterschiedlich weit (ganzer Prozess / Abschnitt / einzelner Schritt)
%  6. Diese Arbeit: drei Ebenen wie Sai et al.; Prozess als strukturierte Liste (1 Prozess, 6 Teilprozesse, 19 Tasks), nicht als BPMN-Grafik
\offen{2.1 -- wird Satz für Satz geschrieben}

\subsection{Regulatorische Compliance}
\offen{2.2}

\subsection{Rechtsrahmen des Anwendungsfalls}
\offen{2.3}

\subsection{Sprachmodelle und Agenten}
\offen{2.4}

\subsection{Bewertungsmaße}
\offen{2.5}
